\section{Liveness and Quit Handlers}
\label{sec:liveness}

Vegas treats the operational question of \emph{what happens when
someone stops responding} as part of the surface syntax rather
than as an implementation detail.  This section walks through the
design: how the language exposes the choice to the programmer,
how the compiler realises it on the EVM, and where the model is
honest about its current limits.

\subsection{Why liveness is in the language}

A multi-party on-chain protocol is exposed to two kinds of
slow-or-absent participants: strategic ones (who decline to play
because the next move is not in their interest), and incapable
ones (client crashes, mempool censorship, gas exhaustion).  The
deployed contract has to decide what to do in both cases, and the
decision must be a total function of public history.  Vegas
makes that decision part of the source: every action that could
plausibly fail to occur carries an explicit handler:

\begin{lstlisting}[style=vg]
yield or split A(x: int);
\end{lstlisting}

The handler describes what the contract does, and the
game-theoretic analysis sees \TT{Quit} as a strategic move
available to the actor.  The two views agree by construction.

\subsection{The three group handlers}

A statement involving one or more roles in a single
\TT{yield}/\TT{commit}/\TT{reveal} group carries one of three
handlers.

\paragraph{split} The quitter's deposit is split among the
surviving roles (in a two-player game, the survivor gets it all).
This is the default for adversarial games: nobody profits by
walking out.

\paragraph{burn} The quitter's deposit is sent to a designated
burn address (or zero address, depending on chain conventions).
Useful in games where giving the deposit to surviving roles would
warp incentives --- voting schemes where a missing voter should
not enrich the others, for instance.

\paragraph{null} The missing field is treated as \TT{null} at the
typing level (the field's type becomes \TT{opt T}), the node's
\TT{done\_} flag is set by the contract, and the protocol
proceeds.  The \TT{withdraw} clause is required to handle the
optional case explicitly.  Useful for protocols where a
non-response is itself strategically meaningful (a delegate
failing to delegate).

\subsection{Per-query payoff handlers}
\label{sec:liveness:per-query}

A single multi-role statement sometimes needs different outcomes
depending on which role quit.  The language supports this via a
\TT{||} qualifier on the role's query:

\begin{lstlisting}[style=vg]
yield or split Bob(ok: bool) || { Alice -> 40 };
\end{lstlisting}

The right-hand side of \TT{||} is an \emph{outcome}: a payoff
allocation chosen from the same sublanguage that \TT{withdraw}
clauses use (a brace-block payoff, or \TT{split}/\TT{burn}/\TT{null}).
When the named role quits, the contract executes the handler's
outcome \emph{instead of} the rest of the protocol; control does
not return.  No subgame or callable game block is involved ---
the handler is a terminal allocation chosen at the bail point.

\subsection{Operational realisation}

The Solidity backend implements \TT{Quit} via a timeout-and-bail
mechanism (\S\ref{sec:solidity:bail}).  The contract maintains one
global inactivity window: if no action has happened anywhere for
\TT{TIMEOUT} seconds (default 24 hours), the next caller's
modifier marks the named role as \TT{bailed} and resets the
window.  Once bailed, the role's actions are rejected by the
\TT{by(Role)} modifier; the bail is one-way.

This single-window design is operationally simple and
implementable in a few lines of Solidity, but its attribution is
structural rather than analytical: it marks ``the role pointed at
by the modifier'' as bailed, not ``the role that was supposed to
act next''.  In a long protocol with several outstanding actions
in different roles, this attribution can be coarser than one would
ideally want.  A per-action deadline scheme would address this at
the cost of substantially more state; the current design is the
simpler starting point.

\subsection{Liveness in the game model}

For the Gambit backend, \TT{Quit} is an additional alphabet
symbol at every move site with a handler.  Equilibrium analyses
therefore see \TT{Quit} as a real move, with the payoff indicated
by the handler.  This is what lets the analysis surface griefing
incentives without modelling them outside the game.  In a
simultaneous \TT{or split} statement, for instance, \TT{Quit} is
strictly dominated by any non-quit move iff the split penalty
exceeds the worst non-quit payoff; the dominance check on the
generated EFG verifies the property mechanically.

\subsection{Limits today}

Two known limits worth flagging.

\paragraph{Fair-play liveness is not statically checked.}  A
\emph{fair-play guard} is a \TT{where} clause whose role can
always find a satisfying value, regardless of prior moves.  The
language does not verify that today; a contextual guard like
\TT{yield B(y: int) where y != A.x \&\& y != A.x + 1 \&\& y != A.x - 1}
over a type \TT{int = \{0..2\}} becomes unsatisfiable depending on
\TT{A}'s play, forcing \TT{B} to quit through no fault of its own.
The project's design notes sketch a static check
(exhaustive enumeration over bounded types, or an SMT call for
the general case) that would type the affected field as
\TT{opt T} when satisfiability cannot be established; implementing
it is on the TODO list under the name \emph{fair-play liveness}.

\paragraph{Computational tractability is not checked either.}  A
guard such as \TT{p * q == N} for a large semiprime \TT{N} is
satisfiable by definition but practically infeasible to satisfy.
Vegas treats this the same as a satisfiable guard; the player
will quit, the contract will bail, and the game model will
record that outcome via the handler.  No analysis distinguishes
``empty move space'' from ``intractable move space''.  This is
the standard limit of purely structural reasoning about action
availability.
